\documentclass[10pt,a4paper]{amsart}
\usepackage[margin=19mm]{geometry}
\usepackage{amsmath,amssymb,mathtools}
\usepackage[scr=rsfs,cal=euler]{mathalfa}
\usepackage{xfrac}
\usepackage[unicode,hidelinks,bookmarks=false]{hyperref}
\newcommand{\ie}{i.e.\ }
\newcommand{\cf}{cf.\ }
\newcommand{\resp}{respectively\ }
\newcommand{\Subsection}{\subsection}
\providecommand{\nobreakdash}{\nobreak\hbox{-}}
\setlength{\emergencystretch}{2em}
\multlinegap0pt
\usepackage{amssymb}
\usepackage{mathtools}
\usepackage[shortlabels]{enumitem}
\usepackage{csquotes}
\usepackage{xcolor}
\usepackage{tikz}
\newtheorem{lemma}{Lemma}
\newtheorem{theorem}{Theorem}
\usepackage{cleveref}
\crefname{lemma}{Lemma}{Lemmas}
\crefname{theorem}{Theorem}{Theorems}
\newcommand{\C}{\mathbb C}
\newcommand{\R}{\mathbb R}
\newcommand{\eps}{\varepsilon}
\title{When Finite Free Curves Split}
\author{Baran Hashemi, Jihoon Hyun}
\date{Source: arXiv:2609.10367v1 (CC BY 4.0)}
\begin{document}
\maketitle

\noindent\emph{Source and licence.} When Finite Free Curves Split. Version arXiv:2609.10367v1 (CC BY 4.0), distributed by arXiv under the Creative Commons Attribution 4.0 International licence, http://creativecommons.org/licenses/by/4.0/, as recorded in the arXiv licence metadata for this exact version and shown on the abstract page https://arxiv.org/abs/2609.10367v1. Full licence notice: \texttt{licenses/arxiv-2609.10367-CC-BY-4.0.txt}.\par\smallskip

\noindent\textit{Editorial scope.} Counted unit: the article's theorem thm:tangent-convolution (source0.tex lines 1594--1631). The article's complete original proof of it is delivered with it (source0.tex lines 1633--1700, one \texttt{proof} environment of 68 lines). No mathematics is written, repaired or re-derived: the statement and the proof are the article's own lines, in the article's own notation, and no retained line is substituted or re-flowed.  Retained uncounted as the source-local context the proof needs: lines 808--817.  Nothing was omitted from any retained span, so the package carries no omission record.  No retained line contains a citation, so the wrapper carries no bibliography and no bibliography entry.  No retained line contains a figure environment, an image-inclusion command or drawing code, so no image is delivered and the package declares no asset dependency.  Editorial formatting only, in addition: an AMS \texttt{amsart} preamble that restates the article's own declarations and macros used by the retained lines, namely the environments \texttt{lemma}, \texttt{theorem} on separate counters, together with the standard packages those lines need.  The retained environments print in this document as Lemma 1, Theorem 1 in order of appearance, whereas the article prints the same items with the numbering of its own full document.  The complete native source of the article as distributed in the arXiv e-print for this exact version is bundled unchanged as \texttt{sources/209905/source0.tex}.
\begin{lemma}[Permutation formula for finite free convolution]
\label{lem:permutation-formula}
For \(\alpha,\beta\in\R^n\),
\begin{equation}\label{eq:hyperbolic-F}
 \mathcal F_n(x,\alpha,\beta)
 :=\frac1{n!}\sum_{\pi\in S_n}
 \prod_{r=1}^n(x-\alpha_r-\beta_{\pi(r)})
 =P_\alpha\boxplus_nP_\beta.
\end{equation}
\end{lemma}

\begin{theorem}[Tangent convolution]\label{thm:tangent-convolution}
Let \(n\geq2\), let \(\alpha,\beta,u,v\in\R^n\), and suppose that
\(a,b\) have multiplicities \(m,k\) as above, with
\(r=m+k-n>0\).  Define the monic degree-\(r\) polynomial
\begin{equation}\label{eq:tangent-convolution-polynomial}
 Q(z)=
 \left(\frac{r!}{m!}U^{(m-r)}(z)\right)
 \boxplus_r
 \left(\frac{r!}{k!}V^{(k-r)}(z)\right)
\end{equation}
and the nonzero real constant
\begin{equation}\label{eq:tangent-convolution-constant}
 c_{\mathrm{tan}}=\frac{m!k!}{n!r!}
       \prod_{i\notin I}(a-\alpha_i)
       \prod_{j\notin J}(b-\beta_j).
\end{equation}
There is a polynomial \(R\in\R[z,\eps]\) such that
\begin{equation}\label{eq:tangent-convolution-expansion}
 h_\eps(a+b+\eps z)
   =\eps^r\bigl(c_{\mathrm{tan}}Q(z)+\eps R(z,\eps)\bigr).
\end{equation}
In particular, \(\eps^{-r}h_\eps(a+b+\eps z)\) converges
coefficientwise, and locally uniformly for \(z\in\C\), to \(c_{\mathrm{tan}}Q(z)\).

The \(r\) output roots approaching \(a+b\), after subtraction of
\(a+b\) and division by \(\eps\), converge as an unordered multiset
to the roots of \(Q\).  If these roots are distinct, say
\(\lambda_1<\cdots<\lambda_r\), the corresponding output roots admit
real analytic branches \(x_1(\eps),\ldots,x_r(\eps)\) through the
parameter value zero,
with
\begin{equation}\label{eq:tangent-convolution-branches}
 x_j(\eps)=a+b+\eps\lambda_j+O(\eps^2),
 \qquad 1\leq j\leq r.
\end{equation}
Within this cluster, the branches are increasingly ordered for positive
\(\eps\) and decreasingly ordered for negative \(\eps\).
\end{theorem}

\begin{proof}
Use the permutation formula of \Cref{lem:permutation-formula}.  For a
permutation \(\pi\in S_n\), let \(d_\pi\) be the number of indices
\(i\in I\) for which \(\pi(i)\in J\).  At most \(n-k\) indices
of \(I\) can be matched outside \(J\), so
\(d_\pi\geq m-(n-k)=r\).  Each such internal match contributes
the factor \(\eps(z-u_i-v_{\pi(i)})\) after the substitution
\(x=a+b+\eps z\).  Every summand is therefore divisible by
\(\eps^r\).  A summand with \(d_\pi>r\) contributes nothing to
the quotient at \(\eps=0\).

Consider a matching with exactly \(r\) internal matches, and let
\(S\subset I\), \(T\subset J\) be the matched subsets, both of
cardinality \(r\).  The remaining \(m-r=n-k\) indices of \(I\)
are matched bijectively to the complement of \(J\), and the complement
of \(I\) is matched bijectively to the remaining \(k-r=n-m\)
indices of \(J\).  There are no matches between the two complements.
At \(\eps=0\), these outside matches contribute
\[
 G=\prod_{i\notin I}(a-\alpha_i)
       \prod_{j\notin J}(b-\beta_j),
\]
independently of their assignment.  For a fixed bijection from \(S\)
to \(T\), the two outside assignments can be chosen in
\((n-k)!(n-m)!\) ways.  Summing over the \(r!\) internal
bijections gives \(r!(P_{u_S}\boxplus_rP_{v_T})(z)\).  Consequently,
the value of the quotient at zero is
\begin{equation}\label{eq:tangent-matching-count}
 \frac{G\,r!(n-k)!(n-m)!}{n!}
 \sum_{\substack{S\subset I,\ |S|=r\\ T\subset J,\ |T|=r}}
       (P_{u_S}\boxplus_rP_{v_T})(z).
\end{equation}

We have averaged over the bijections between the selected roots. It remains to average over the choices of the \(r\)-element subsets themselves. The average of the corresponding monic products is a normalized derivative of the original velocity polynomial. Differentiating a product gives
\[
 \frac{r!}{m!}U^{(m-r)}
   =\binom mr^{-1}\sum_{\substack{S\subset I\\|S|=r}}P_{u_S},
 \qquad
 \frac{r!}{k!}V^{(k-r)}
   =\binom kr^{-1}\sum_{\substack{T\subset J\\|T|=r}}P_{v_T}.
\]
The coefficient formula makes convolution bilinear in its polynomial arguments. Hence the double sum in \eqref{eq:tangent-matching-count} is \(\binom mr\binom kr Q\). Its scalar coefficient simplifies to
\(Gm!k!/(n!r!)=c_{\mathrm{tan}}\), because \(n-k=m-r\) and \(n-m=k-r\). This proves \eqref{eq:tangent-convolution-expansion} and the polynomial divisibility also proves both stated forms of convergence.

Fujie's atom theorem says that \(a+b\) has multiplicity exactly \(r\)
in \(h_0\).  The constant can also be checked directly from the
expansion: the coefficient of \(\eps^rz^r\) is
\(h_0^{(r)}(a+b)/r!=c_{\mathrm{tan}}\). The two normalized derivatives in
\eqref{eq:tangent-convolution-polynomial} are real-rooted by Rolle's
theorem, and their convolution is real-rooted by Walsh's theorem.

We have identified the limiting polynomial. We must still show that its roots describe precisely the output cluster approaching \(a+b\).  Choose a circle containing every root of \(Q\), with no root on its boundary. Local uniform convergence and Rouch\'e's theorem show that,
for sufficiently small nonzero \(\eps\), the rescaled polynomial
has exactly \(r\) roots inside this circle.  Their unscaled roots tend
to \(a+b\), so they are precisely its cluster, since \(h_0\) has
multiplicity \(r\) there.  Applying the same argument to disjoint
small circles around the distinct roots of \(Q\) proves multiset
convergence, including their multiplicities.

If \(\lambda_j\) is simple, apply the analytic implicit-function
theorem to \(c_{\mathrm{tan}}Q(z)+\eps R(z,\eps)\) at \((\lambda_j,0)\).
It gives a real analytic solution \(z_j(\eps)\) with
\(z_j(0)=\lambda_j\).  Setting
\(x_j(\eps)=a+b+\eps z_j(\eps)\) proves
\eqref{eq:tangent-convolution-branches}.  The signs of the differences
\(x_j(\eps)-x_i(\eps)=\eps(\lambda_j-\lambda_i)+O(\eps^2)\)
give the assertion about ordering.
\end{proof}

\end{document}
